\section{Methodology}
\label{sec:methodology}

This section formalizes the ranking problem, describes the public benchmark and evaluation harness, and details the autoresearch loop and baselines.

\subsection{Problem Formulation}
\label{sec:problem}

Given a pool of $N$ peptide candidates, each characterized by estimated scores for $K$ endpoints (activity, toxicity, stability, manufacturability), the goal is to learn a ranking policy $\pi$ that maps score vectors to a candidate ordering maximizing discovery utility under all constraints simultaneously.
Formally, let $\mathbf{x}_i = (a_i, t_i, s_i, d_i) \in \mathbb{R}^4$ denote the endpoint score vector for candidate $i$.
A ranking policy $\pi: \mathbb{R}^{N \times 4} \to \mathcal{S}_N$ produces a permutation of candidates.
The policy is evaluated by how well its top-$k$ selections overlap with ground-truth high-quality candidates defined by independent oracle functions (Section~\ref{sec:evaluation_harness}).

This formulation differs from generative molecular optimization, which proposes novel molecules.
Here, the candidate pool is fixed and the challenge is \emph{prioritization}: given imperfect multi-endpoint estimates, which candidates should advance to expensive downstream validation?

\subsection{Public Benchmark}
\label{sec:benchmark}

We assemble a benchmark from four public data sources spanning complementary developability endpoints (Table~\ref{tab:datasets}).

\begin{table}[t]
\centering
\caption{Benchmark data sources. The candidate pool comprises 3{,}554 antimicrobial peptides from DBAASP with measured activity against \emph{E.~coli} ATCC 25922. Toxicity and stability are predicted by random forest models trained on independent public datasets. Developability is a deterministic rule-based penalty. Sequence lengths range from 5 to 50 residues.}
\label{tab:datasets}
\small
\begin{tabular}{@{}llrrll@{}}
\toprule
\textbf{Endpoint} & \textbf{Source} & \textbf{$N$} & \textbf{Type} & \textbf{Held-out} & \textbf{Notes} \\
\midrule
Activity & DBAASP v4 & 3{,}554 & Measured MIC & --- & $-\log_{10}$ potency \\
Toxicity & ToxinPred3 & 8{,}323 train & RF classifier & AUC 0.920 & 8.9\% overlap flagged \\
Stability & HLP & 1{,}353 train & RF regressor & $R^2$ 0.547 & Weak; acknowledged \\
Developability & Rule-based & --- & 7 penalties & Deterministic & Length, charge, hydrophobicity \\
\bottomrule
\end{tabular}
\end{table}

\textbf{Activity.}
Antimicrobial activity is derived from minimum inhibitory concentration (MIC) values in the Database of Antimicrobial Activity and Structure of Peptides (DBAASP v4) \citep{pirtskhalava2021dbaasp}, filtered to \emph{E.~coli} ATCC 25922 as target strain.
Raw MIC values are converted to $-\log_{10}$ potency and min-max normalized to $[0, 1]$.
This is the only ground-truth endpoint; the remaining three are estimated.

\textbf{Toxicity.}
A random forest classifier is trained on binary toxicity labels from ToxinPred3 \citep{rathore2024toxinpred3}, using physicochemical features (amino acid composition, hydrophobicity, charge, molecular weight) augmented with ESM-2 embeddings \citep{lin2023esm2} from the 8M-parameter model.
Held-out AUC is 0.920.
An audit identified 8.9\% sequence overlap between the ToxinPred3 training set and DBAASP candidate pool; this overlap is flagged but not fully controlled in the current benchmark.

\textbf{Stability.}
A random forest regressor is trained on peptide half-life data from HLP \citep{sharma2014hlp}, using the same physicochemical feature set.
Held-out $R^2$ is 0.547, substantially weaker than the toxicity model; this limitation propagates to any ranking policy that relies on stability scores.

\textbf{Developability.}
A deterministic rule-based proxy computes a penalty score from seven sequence-level rules: length outside 8--40 residues ($+1$), consecutive hydrophobic stretches ($+0.5$ per excess), consecutive prolines ($+0.5$), high net charge ($+1$), high mean hydrophobicity ($+1$), and excess cysteines ($+1$).
The total penalty ranges from 0 (ideal) to ${\sim}3.5$ in practice.

\textbf{Candidate pool statistics.}
The pool of 3{,}554 peptides has mean activity $0.418 \pm 0.124$, mean toxicity $0.662 \pm 0.163$, mean stability $0.205 \pm 0.108$, and mean developability penalty $0.780 \pm 0.661$.
The data are split approximately 70/15/15\% into training, validation, and test sets (2{,}452/539/563 candidates).
All results in this paper are reported on the validation split unless otherwise noted.

\subsection{Fixed Evaluation Harness}
\label{sec:evaluation_harness}

A fixed evaluation harness prevents the agent from gaming any single metric definition.
The harness defines three independent oracle functions, each representing a different notion of ``ground-truth best'' candidates:

\textbf{Pareto rank oracle.}
Each candidate is assigned its Pareto domination count across all four endpoints (treating activity and stability as objectives to maximize, toxicity and developability penalty as objectives to minimize).
Candidates with fewer dominating solutions rank higher.
This oracle is resistant to linear approximation.

\textbf{Rank product oracle.}
The geometric mean of per-endpoint ranks is computed: $\text{score}_i = (r^\text{act}_i \cdot r^\text{tox}_i \cdot r^\text{stab}_i \cdot r^\text{dev}_i)^{1/4}$.
This nonlinear aggregation penalizes candidates that are weak on any single endpoint.

\textbf{Threshold-gated oracle.}
Activity is multiplied by a stability bonus and gated by sigmoid functions on toxicity and developability: $\text{score}_i = a_i \cdot (0.3 + 0.7 \cdot s_i) \cdot \sigma(-10(t_i - 0.5)) \cdot \sigma(-10(d_i - 0.5))$.
This oracle sharply penalizes candidates exceeding toxicity or developability thresholds.

\textbf{Metrics.}
All metrics are computed as the mean across all three oracles:
\begin{itemize}[nosep]
  \item \textbf{Top-$k$ enrichment} (primary): fraction of the oracle's true top-$k$ appearing in the policy's top-$k$ ($k = 20$).
  \item \textbf{NDCG} \citep{jarvelin2002ndcg} (secondary): normalized discounted cumulative gain at $k = 20$.
  \item \textbf{Hypervolume} \citep{zitzler2003hypervolume} (diagnostic): dominated area in the (activity, $1 - \text{toxicity}$) plane for the policy's top-$k$ selections.
  \item \textbf{Feasible fraction} (constraint): fraction of top-$k$ selections satisfying $\text{toxicity} < 0.5$ and $\text{dev\_penalty} \leq 3$.
\end{itemize}

Using three oracles with different mathematical structures (Pareto counting, geometric aggregation, and threshold gating) prevents the agent from reverse-engineering any single formula.

\subsection{Autoresearch Loop}
\label{sec:loop}

We adapt the autoresearch framework \citep{karpathy2026autoresearch} from neural architecture search to ranking policy search.
Our companion paper \citep{wijaya2026autonomousagentdiscoversmolecular} previously applied this paradigm to molecular transformer architecture search across SMILES, protein, and NLP domains (3{,}106 experiments).
Here, the editable surface changes from a training script to a ranking policy, and the evaluation changes from a single metric (val\_bpb) to a multi-objective ensemble.

\textbf{Architecture.}
An outer loop (\texttt{run\_agent\_loop.py}) spawns one-shot Codex sessions (GPT-5.4, xhigh reasoning effort).
Each session receives: (1)~the system prompt (\texttt{program.md}) specifying loop rules, (2)~the current ranking policy (\texttt{rank.py}), and (3)~the results log of all prior experiments.
The session edits \texttt{rank.py} to implement one coherent policy modification, then invokes the evaluation harness.
If the primary metric (mean top-$k$ enrichment) improves, the change is kept; otherwise it is reverted.

\textbf{Workspace isolation.}
Each session operates on an isolated copy of the source code.
The evaluation harness, data pipeline, endpoint models, and oracle definitions are frozen throughout the loop.
The agent can only modify \texttt{rank.py} (and optionally \texttt{rank\_learnable.py} for learned models).

\textbf{Budget.}
The loop executes 100 sequential experiments.
Of these, 12 are kept (11 improvements plus the initial baseline) and 88 are discarded (regressions or ties).

\subsection{Baselines}
\label{sec:baselines}

We compare the agent-improved policy against seven baselines spanning simple heuristics, standard multi-objective optimization, and random search:

\begin{itemize}[nosep]
  \item \textbf{Activity only}: rank by predicted activity alone.
  \item \textbf{Toxicity exclusion}: hard-filter candidates with toxicity $\geq 0.5$, rank remainder by activity.
  \item \textbf{Weighted sum}: min-max normalize all four endpoints, combine with equal weights ($w = 0.25$ each).
  \item \textbf{Rule only}: hard-filter on toxicity $\leq 0.7$ and developability penalty $\leq 3$, rank by stability then activity.
  \item \textbf{Random}: uniformly random ordering (seed = 42).
  \item \textbf{Random weight search}: best-of-1{,}000 Dirichlet-sampled weight vectors evaluated with the weighted-sum scalarization \citep{bergstra2012random}. This controls for blind exploration of the same policy class the agent starts from.
  \item \textbf{NSGA-II} \citep{deb2002nsga2}: non-dominated sorting with crowding distance tie-break across all four objectives. This is the standard multi-objective optimization baseline.
\end{itemize}
